\subsection{整元素的刻画}

定理 3. 设 A 是域 K 的子环。A 在 K 中的整闭包 A' 是 K 的包含 A 的诸赋值环之交；若 A 是局部的，则 A' 是 K 的支配 A 的诸赋值环之交。

设 \(x\) 是 \(A'\) 的元素，\(V\) 是 \(K\) 的包含 \(A\) 的赋值环；由于 \(x\) 在 \(V\) 上整，存在素理想 \(p'\)（\(V[x]\) 的）使 \(p' \cap V = m(V)\)（第五章 §2 第 1 节定理 1）；显然局部环 \((V[x])_{p'}\) 支配 \(V\)，从而等于它；于是 \(x \in V\)。反之，设 \(y\) 是 \(K\) 的不在 \(A\) 上整的元素；则存在极大理想 \(M\)（\(A[y^{-1}]\) 的），它包含 \(y^{-1}\)（第 2 节引理 1）；还存在赋值环 \(V\)（\(K\) 的），它支配 \((A[y^{-1}])_{m}\)（第 2 节定理 2 的推论）；由于 \(y^{-1} \in m(V)\)，有 \(y \notin V\)。此外 \(M \cap A\) 是 \(A\) 的极大理想（第 2 节引理 1）；因此，若 \(A\) 是局部的，则 \(M \cap A = m(A)\)，从而 \(V\) 支配 \(A\)。

推论 1. 每个赋值环都是整闭的。

推论 2. 为使一个整环是整闭的，必须且只需它是它的分式域的一族赋值环之交。

在 Noether 环的情形，推论 2 可以更精确化（第七章 § 1 第 3 节定理 1 的推论）。

推论 3. 设 K 是域，K' 是 K 的扩张，A 是 K 的赋值环。A 在 K' 中的整闭包是 K' 的满足 V' ∩ K = A 的诸赋值环 V' 之交。

定理 1 (c) 表明，若 \(V'\) 是 \(K'\) 的赋值环，则 \(V' \cap K\) 是 K 的赋值环，且 \(V'\) 支配 \(V' \cap K\)。为使 \(V'\) 支配 A，必须且只需 \(V' \cap K\) 支配 A，从而等于 A。

